\documentclass[12pt]{article}
\usepackage[margin=1in]{geometry}
\usepackage{amsmath,amssymb,amsfonts}
\usepackage{graphicx}
\usepackage{booktabs}
\usepackage{hyperref}
\usepackage{authblk}
\usepackage{setspace}
\usepackage{xcolor}
\onehalfspacing

\newcommand{\si}[1]{\mathrm{#1}}
\newcommand{\TODO}[1]{{\color{red}\textbf{[#1]}}}

\title{Sample-Efficient MAPPO for Physics-in-the-Loop Multi-Robot Training:\\
Diagnosing and Correcting Entropy Collapse Under a\\
Thousand-Step, Not Million-Step, Budget}
\author[1]{Harsh Pandhe}
\affil[1]{Independent Research \\ \texttt{harshpandhehome@gmail.com}}
\date{}

\begin{document}
\maketitle

\begin{abstract}
\noindent
Multi-agent reinforcement learning (MARL) sample-efficiency work
overwhelmingly assumes access to a fast, often GPU-parallelized simulator,
where the practical constraint is wall-clock per environment step, not the
number of steps a single rollout worker can collect. This assumption breaks
down for physics-in-the-loop robotics training against a general-purpose
simulator (Gazebo Sim) with no massively-parallel rollout path: each
environment step costs real wall-clock time regardless of parallelization
strategy, and a training budget of thousands, not millions, of environment
steps is a hard practical ceiling, not a simplifying choice. We report a
directly observed failure mode under this constraint -- a centralized-critic
MAPPO policy for cooperative multi-robot area coverage, trained for
$4{,}500$ total environment steps ($15$ iterations, train batch size $300$),
whose policy entropy collapsed to $-0.53$ by the final iteration while
episode returns swung between $-227$ and $+3$ within the same $3$--$4$
sampled episodes -- and diagnose it as a direct consequence of leaving
RLlib's advantage-estimation ($\lambda=1.0$) and exploration-pressure
($\text{entropy\_coeff}=0.0$) hyperparameters at defaults tuned for
much larger batch regimes. We derive a specific, mechanistically justified
correction (reduced GAE $\lambda$, a nonzero entropy bonus, reward
rebalancing to reduce per-episode return variance, and batch/episode-length
adjustments that increase distinct-episode count per iteration without
increasing wall-clock), implement it in a public codebase, and report the
diagnosis, the fix, and its expected mechanism in full; we could not
complete the corrected training run's full empirical validation in this
revision due to an unrelated infrastructure fault (a segmentation fault in
the local PyTorch/CUDA installation, confirmed independent of the training
code by three reproducible crash attempts) and report this transparently as
an open item alongside the diagnosis and proposed fix.
\end{abstract}

\noindent\textbf{Keywords:} multi-agent reinforcement learning, sample
efficiency, MAPPO, proximal policy optimization, robotics, entropy collapse

\section{Introduction}
Sample efficiency in deep reinforcement learning is usually studied under an
implicit assumption: environment steps are cheap enough, or parallelizable
enough, that the real constraint is algorithmic (how much a policy improves
per gradient update) rather than infrastructural (how many environment steps
can physically be collected before an experiment's wall-clock budget is
spent). GPU-accelerated physics simulators built for exactly this
assumption -- running thousands of environment instances
concurrently -- have become the default substrate for sample-efficiency
research in robot learning. This is a reasonable choice when the target
deployment simulator supports it. It is not a universal property of robotics
simulation: general-purpose robotics simulators used for realistic sensor
modeling, ROS integration, and physically detailed multi-robot interaction
(collision geometry, LiDAR ray-casting, differential-drive dynamics) are
frequently run as a single real-time (or near-real-time) instance, because
the fidelity that makes them useful for the target application is exactly
what makes naive massive parallelization impractical without dedicated
infrastructure investment most individual researchers and small teams do not
have. Under this constraint, a training budget of thousands of environment
steps -- not the millions typical of parallelized-simulator MARL
benchmarks -- is not a choice, it is the wall-clock-bounded ceiling.

We report a concrete instance of what happens to a standard MARL algorithm
(MAPPO~\cite{yu2022mappo} with a centralized critic, CTDE) under exactly this
constraint, diagnose the failure mechanistically rather than empirically
alone, and derive a correction from that mechanism. Our contribution is not
a new algorithm; it is a grounded case study connecting a specific,
reproducible failure (entropy collapse under a tiny, wall-clock-bounded
sample budget) to specific hyperparameter and reward-design choices, in a
setting -- physics-in-the-loop multi-robot training with no fast-simulator
escape hatch -- that the dominant sample-efficiency literature does not
directly address.

\section{Related Work}
Yu et al.~\cite{yu2022mappo} establish MAPPO with a centralized critic as
a strong, simple CTDE baseline for cooperative multi-agent games, and is the
algorithm family we build on; their sample-efficiency results, like most of
the literature that follows, are reported in a regime with orders of
magnitude more environment steps than are available under our constraint.
Sample-efficiency techniques aimed explicitly at closing this gap
increasingly rely on differentiable simulation or model-based rollouts to
reduce real-environment sample requirements by large
factors~\cite{asymmetricphysics2026,agilemotor2024} -- a complementary but
architecturally distinct strategy from the hyperparameter- and
reward-design-level correction we study, and one that requires a
differentiable or learned dynamics model not assumed available here. Masked
reconstruction and other representation-learning approaches improve MARL
sample efficiency by extracting more training signal per environment
step~\cite{maskedreconstruction2023}, a strategy compatible with and
potentially complementary to the correction we propose, though we do not
implement it here. Physics-informed MARL for distributed multi-robot
problems~\cite{physicsinformedmarl2024} incorporates physical structure
directly into the learning problem, again a complementary direction. None of
this literature, to our knowledge, specifically characterizes the entropy
collapse and advantage-variance failure mode we report as a direct
consequence of applying default single-agent PPO hyperparameters
unmodified in a small-batch, physics-in-the-loop multi-agent setting.

\section{System and Training Setup}
\label{sec:system}
The learned system is a shared MAPPO policy for a three-robot TurtleBot3
swarm performing cooperative area coverage in Gazebo Sim (ROS~2 Jazzy),
trained via Ray RLlib. The actor is a permutation-invariant network:
ego features ($28$-dim: LiDAR, goal, velocity) are concatenated with a
masked mean-pooled embedding of up to $9$ neighbors' relative
$(\text{distance}, \text{angle})$, each processed by a shared two-layer MLP.
The critic is centralized: it additionally consumes a masked mean-pooled
embedding of all teammates' joint (observation, action) pairs, following the
standard CTDE recipe. Training uses \texttt{num\_env\_runners=0} (Gazebo must
run in-process; RLlib's usual multi-worker parallel-rollout path is
unavailable), so every environment step is collected serially against
real-time (or near-real-time) physics -- the exact constraint motivating this
paper.

\textbf{Baseline configuration (the diagnosed failure).} $15$ training
iterations, train batch size $300$, minibatch size $64$, $5$ epochs per
iteration, learning rate $1\times10^{-4}$, PPO clip $0.2$, discount
$\gamma=0.99$, episode length (\texttt{max\_steps}) $150$,
\texttt{rollout\_fragment\_length} $150$. Critically, GAE $\lambda$ and
\texttt{entropy\_coeff} were left at RLlib's PPO defaults ($1.0$ and $0.0$
respectively) -- i.e., no advantage-variance reduction and no exploration
bonus. Total training budget: $15 \times 300 = 4{,}500$ environment steps.

\section{Diagnosis}
\label{sec:diagnosis}
With train batch size $300$ split across $3$ agents over episodes up to
$150$ steps, a single training iteration collects on the order of
$3$--$4$ complete episodes. Two consequences follow directly from the
baseline configuration:

\textbf{(1) High-variance advantage estimates.} GAE $\lambda=1.0$ corresponds
to the (unbiased but high-variance) Monte-Carlo return estimator. With only
$3$--$4$ episodes contributing to a batch, the resulting advantage estimates
are dominated by whichever few episodes happened to be sampled, rather than
reflecting the policy's average behavior. We observed this directly: episode
returns within a single iteration's batch ranged from $-227$ to $+3$ --
a swing far larger than plausible in-policy variation, indicative of an
estimator amplifying a small sample rather than measuring a stable signal.

\textbf{(2) No counter-pressure against premature convergence.}
\texttt{entropy\_coeff}$=0.0$ means the PPO objective contains no term
rewarding continued exploration. Combined with (1) -- a high-variance,
noisy advantage signal -- the policy has every incentive to converge quickly
toward whatever locally low-variance (e.g.\ near-stationary, ``safe'')
behavior it stumbles into early, since that behavior at least yields a
consistent, low-magnitude return, rather than continuing to explore into the
noisy, sparsely-rewarded regions of the state space where the actual task
reward (area coverage) lives. We observed this directly: policy entropy fell
to $-0.53$ by the final (15th) iteration, having been positive earlier in
training -- a collapse toward a low-entropy, near-deterministic policy well
before $4{,}500$ environment steps could plausibly be sufficient to have
converged on a genuinely good policy for a $3$-robot coordination task.

A third, compounding factor is reward design: the per-step reward combines
a small dense progress/coverage term with sparse $\pm20$ spikes for
goal-reaching and collision. Over an episode as short as $150$ steps
(or fewer, if the episode terminates early on collision), a single $-20$
collision event can dominate the cumulative return, further inflating the
return variance described in (1) independent of the $\lambda$ setting alone.

\section{Proposed Correction}
\label{sec:correction}
We derive a correction directly from the mechanisms in
Section~\ref{sec:diagnosis}, rather than a general hyperparameter sweep:

\begin{itemize}
  \item \textbf{GAE $\lambda: 1.0 \to 0.95$.} Directly targets mechanism (1):
        trades a small amount of estimator bias for a substantial reduction
        in advantage-estimate variance, which matters more, not less, when
        the batch contains only a handful of episodes.
  \item \textbf{\texttt{entropy\_coeff}: $0.0 \to 0.01$.} Directly targets
        mechanism (2): introduces continued exploration pressure
        proportional to policy entropy, intended to prevent the observed
        collapse to $-0.53$ before the policy has had a chance to discover
        the coverage-reward signal.
  \item \textbf{Reward rebalancing.} Collision penalty $-20.0 \to -8.0$
        (reduces the dominance of a single sparse spike over a short
        episode); coverage-reward weight $2.0 \to 4.0$ (densifies the signal
        for the metric actually evaluated, partially independent of
        mechanism (1)--(2) but compounding with them by reducing overall
        return variance).
  \item \textbf{Episode length and batch size: \texttt{max\_steps}
        $150\to90$, train batch size $300\to700$.} Episode-reset overhead in
        a physics-in-the-loop environment (settling time after simulator
        reset) is fixed regardless of episode length, so shorter episodes
        yield more distinct episodes per unit wall-clock; a larger batch
        size directly increases the number of episodes contributing to each
        gradient step, compounding with the $\lambda=0.95$ variance
        reduction. Both changes increase distinct-episode count per
        iteration without a proportional wall-clock increase, respecting
        the paper's central constraint (Section 1) that environment steps
        cannot simply be made cheaper.
  \item \textbf{Iteration budget: $15\to45$.} A secondary lever, applied only
        after the above changes, within an estimated wall-clock budget of
        $1.6$--$3.5$ hours based on observed per-iteration timing at the
        new batch/episode configuration.
\end{itemize}

We explicitly considered and rejected two further interventions given the
project's wall-clock budget. \emph{Curriculum learning} (training on $1$
robot before scaling to $3$) would require reworking spawn and goal logic
for a variable robot count and a staged training driver, at a cost of
multiple engineering hours, to address a credit-assignment difficulty that
the centralized critic already handles structurally -- we judged this a
speculative fix for a problem not diagnosed in Section~\ref{sec:diagnosis}.
\emph{Behavior-cloning warm start} from a working non-learned frontier
heuristic (available in the same codebase) would require an offline dataset
pipeline, a supervised pretraining loop, and reconciling the heuristic's
discrete waypoint targets with the continuous action space and the
centralized critic's value function -- a distinct problem (cold-start
exploration) from the one diagnosed here (entropy collapse from
advantage-variance under a tiny batch), and correspondingly not pursued.

\section{Validation Status}
\label{sec:validation}
The correction in Section~\ref{sec:correction} is implemented in the
project's training script and verified not to regress existing
functionality: the unit test suite (including a centralized-critic
postprocessing test and a permutation-invariance test at swarm sizes unseen
during training) passes unchanged, and a companion coverage-demo pipeline
sharing the same environment code was independently re-verified end-to-end
after the change. \TODO{We were unable to complete the corrected
configuration's full training run in this revision: three independent
launch attempts (including one after disabling PyTorch's dynamo/AOT-compile
subsystem, the standard mitigation for the specific crash signature
observed) terminated with an identical, reproducible native segmentation
fault inside \texttt{torch.optim.adam.Adam.\_\_init\_\_}, confirmed by
matching stack traces across all three attempts and unrelated to any
training-code or hyperparameter change -- i.e., a local PyTorch/CUDA
installation fault, not a finding about the correction itself. We report
this transparently rather than fabricate or omit the pending result. A
revision of this paper will report post-correction entropy trajectory,
episode-return variance, and downstream task-coverage numbers once training
infrastructure is repaired; the diagnosis (Section~\ref{sec:diagnosis}) and
correction rationale (Section~\ref{sec:correction}) do not depend on that
outcome and are reported in full here.}

We note, as a methodological point relevant to this paper's own thesis, that
this infrastructure fault is itself an instance of the operational fragility
that distinguishes physics-in-the-loop training from cloud-scale,
parallelized-simulator MARL research: a single-machine, single-instance
training pipeline has a single point of failure that a fleet of
identical parallel simulator instances does not, and debugging it competes
directly for the same limited wall-clock budget that motivates this paper's
sample-efficiency concern in the first place.

\section{Discussion}
The failure mode diagnosed here -- entropy collapse from high-variance
advantage estimates under a batch too small for $\lambda=1.0$, compounded by
a lack of exploration pressure and a spiky reward -- is not specific to
multi-robot coverage or to Gazebo; it is a generic consequence of applying
PPO/MAPPO default hyperparameters, tuned in a literature that assumes
large-batch training, to any setting where the environment-step budget is
externally bounded by real-time physics rather than by researcher choice.
We expect the same diagnosis (if not the identical numeric correction) to
transfer to other physics-in-the-loop, single-instance-simulator MARL
settings: any setup where \texttt{num\_env\_runners} cannot exceed a small
number for infrastructural reasons is exposed to the same batch-size/GAE
interaction. We see the primary value of this paper as making the diagnosis
and correction explicit and mechanistically grounded, rather than leaving
practitioners to rediscover it empirically, run by expensive run, under a
sample budget that makes empirical hyperparameter search itself costly.

\section{Conclusion}
Under a hard, wall-clock-bounded sample budget of thousands rather than
millions of environment steps -- the realistic regime for physics-in-the-loop
multi-robot training without dedicated parallel-simulation infrastructure --
default PPO/MAPPO hyperparameters tuned for much larger batch regimes
produce a specific, reproducible failure: high-variance advantage estimates
from $\lambda=1.0$, compounded by zero exploration pressure, collapsing
policy entropy before the policy has seen enough data to learn the task. We
diagnosed this failure mechanistically in a real multi-robot coverage
system, derived a targeted correction (reduced $\lambda$, nonzero entropy
bonus, reward rebalancing, and episode/batch sizing that increases
distinct-episode count without increasing wall-clock), and report both the
diagnosis and the correction in full. Completing the correction's empirical
validation was blocked in this revision by an infrastructure fault unrelated
to the training methodology, reported transparently rather than concealed;
we consider the mechanistic diagnosis and grounded correction, independently
useful to practitioners facing the same constraint, to stand on their own
pending that validation.

\bibliographystyle{plain}
\bibliography{topic1_references}

\end{document}
